%% ma: L12-B12-0025
%% lop: 12
%% chuong: 4
%% bai: 12
%% dang: 12.12.17
%% loai: DS
%% muc: [NB, TH, VD, VD]
%% dap_an: "ĐĐSS"
%% nguon: "(NBV)-ĐỀ SỐ 19-MỖI NGÀY 1 ĐỀ THI 2025.pdf (D19, MỖI NGÀY 1 ĐỀ - Đề số 19), Phần II Câu 3"
%% kien_thuc: "Bài 12 (ứng dụng tích phân: tổng doanh thu, tổng chi phí từ tốc độ thay đổi), Bài 5 (đạo hàm, lợi nhuận biên)"
%% trang_thai: cho-duyet
%% ly_do: "Dạng mới đề xuất 12.12.17 (doanh thu và chi phí của máy công nghiệp theo tốc độ biên: tích phân, tuổi thọ hữu ích, lợi nhuận), chờ giáo viên duyệt."
%% ghi_chu: "Đáp án gốc Đ Đ S S khớp. Lời giải gốc làm tròn tuổi thọ hữu ích thành 6 năm; thực ra T = √(602/17) ≈ 5,95 năm và lợi nhuận ≈ 2388,3 triệu đồng (giải độc lập), kết luận c), d) vẫn là sai. Bỏ ảnh minh họa chiếc máy (ảnh trang trí)."
\begin{ex}
Giả sử rằng khi được $t$ năm tuổi, một máy công nghiệp $A$ tạo ra doanh thu với tốc độ $R'(t)=650-4t^2$ (triệu đồng/năm), thời điểm $t=0$ tính từ lúc máy $A$ bắt đầu hoạt động. Biết rằng chi phí biên cho vận hành và bảo trì là $C'(t)=48+13t^2$ (triệu đồng/năm), ở đây $C(t)$ là chi phí vận hành và bảo trì của máy $A$ khi nó được $t$ năm tuổi.
\choiceTF
{\True Doanh thu sau $10$ năm của máy $A$ là $\displaystyle\int_0^{10}\left(650-4t^2\right)\mathrm dt$ (triệu đồng).}
{\True Tổng chi phí vận hành và bảo trì của máy $A$ trong $6$ năm là $1224$ (triệu đồng).}
{Tuổi thọ hữu ích của một máy là số năm $T$ trước khi lợi nhuận (bằng doanh thu trừ chi phí) mà nó tạo ra bắt đầu giảm. Tuổi thọ hữu ích của máy $A$ này là $7$ năm.}
{Lợi nhuận do máy $A$ tạo ra trong suốt thời gian tuổi thọ hữu ích của nó là $2440$ (triệu đồng).}
\loigiai{a) Doanh thu sau $10$ năm là $\displaystyle\int_0^{10}R'(t)\,\mathrm dt=\int_0^{10}\left(650-4t^2\right)\mathrm dt$ nên a) đúng.

b) Tổng chi phí trong $6$ năm là $\displaystyle\int_0^6\left(48+13t^2\right)\mathrm dt=\left(48t+\dfrac{13t^3}3\right)\Big|_0^6=288+936=1224$ (triệu đồng): b) đúng.

c) Lợi nhuận biên là $P'(t)=R'(t)-C'(t)=602-17t^2$. Lợi nhuận bắt đầu giảm khi $P'(t)<0\Leftrightarrow t>\sqrt{\dfrac{602}{17}}\approx5{,}95$. Vậy tuổi thọ hữu ích khoảng $5{,}95$ năm (xấp xỉ $6$ năm), không phải $7$ năm: c) sai.

d) Lợi nhuận trong thời gian đó là $\displaystyle\int_0^T\left(602-17t^2\right)\mathrm dt=602T-\dfrac{17T^3}3\approx2388$ (triệu đồng), khác $2440$: d) sai.}
\end{ex}
